\documentclass[12pt,a4paper]{article}

%% ── Codificación y fuentes ────────────────────────────────────────
\usepackage[utf8]{inputenc}
\usepackage[T1]{fontenc}
\usepackage{lmodern}
\usepackage[spanish]{babel}

%% ── Página ────────────────────────────────────────────────────────
\usepackage[margin=2cm, top=2.5cm, bottom=2.5cm, headheight=25pt]{geometry}

%% ── Tipografía y colores ──────────────────────────────────────────
\usepackage{xcolor}
\usepackage{enumitem}
\usepackage{array}
\usepackage{booktabs}
\usepackage{colortbl}
\usepackage{tabularx}
\usepackage{multirow}
\usepackage{amsmath}

%% ── Cabecera/pie ──────────────────────────────────────────────────
\usepackage{fancyhdr}
\usepackage{titlesec}

%% ── Cajas ─────────────────────────────────────────────────────────
\usepackage{tcolorbox}
\tcbuselibrary{skins, breakable}

%% ── Iconos ────────────────────────────────────────────────────────
\usepackage{fontawesome5}

%% ── Diagramas ─────────────────────────────────────────────────────
\usepackage{tikz}
\usetikzlibrary{arrows.meta, positioning, calc, shapes.geometric, matrix, fit, backgrounds, babel}

%% ══════════════════════════════════════════════════════════════════
%% COLORES
%% ══════════════════════════════════════════════════════════════════
\definecolor{accessred}{RGB}{175,0,0}
\definecolor{darkblue}{RGB}{0,70,127}
\definecolor{lightblue}{RGB}{210,230,255}
\definecolor{darkgreen}{RGB}{0,110,0}
\definecolor{lightgreen}{RGB}{220,250,220}
\definecolor{chalkwhite}{RGB}{248,248,242}
\definecolor{chalk}{RGB}{238,238,232}
\definecolor{chalkyellow}{RGB}{255,252,180}
\definecolor{chalkred}{RGB}{255,210,210}
\definecolor{board}{RGB}{30,80,50}
\definecolor{timercolor}{RGB}{200,70,0}
\definecolor{actcolor}{RGB}{0,120,170}
\definecolor{theorycolor}{RGB}{90,0,110}
\definecolor{practcolor}{RGB}{0,110,55}
\definecolor{tipcolor}{RGB}{160,80,0}
\definecolor{warncolor}{RGB}{150,0,0}
\definecolor{graylight}{RGB}{248,248,248}

%% ══════════════════════════════════════════════════════════════════
%% CABECERA
%% ══════════════════════════════════════════════════════════════════
\pagestyle{fancy}
\fancyhf{}
\fancyhead[L]{\textcolor{accessred}{\textbf{Guión docente — Sesiones 1 y 2}}}
\fancyhead[R]{\textcolor{gray}{\small UT6: Microsoft Access · AOF}}
\fancyfoot[C]{\textcolor{gray}{\small\thepage}}
\renewcommand{\headrulewidth}{1.5pt}
\renewcommand{\headrule}{%
  \hbox to\headwidth{\color{accessred}\leaders\hrule height\headrulewidth\hfill}}

%% ══════════════════════════════════════════════════════════════════
%% SECCIONES
%% ══════════════════════════════════════════════════════════════════
\titleformat{\section}{\LARGE\bfseries\color{accessred}}{}{0em}{}[\color{accessred}\titlerule]
\titleformat{\subsection}{\large\bfseries\color{darkblue}}{}{0em}{}

%% ══════════════════════════════════════════════════════════════════
%% CAJAS TCOLORBOX
%% ══════════════════════════════════════════════════════════════════
\newtcolorbox{timebox}[1]{
  colback=timercolor!10, colframe=timercolor,
  fonttitle=\bfseries\color{white}, colbacktitle=timercolor,
  title={#1}, left=6pt, right=6pt, breakable, enhanced}

\newtcolorbox{actbox}[1]{
  colback=actcolor!8, colframe=actcolor,
  fonttitle=\bfseries\color{white}, colbacktitle=actcolor,
  title={#1}, breakable, enhanced}

\newtcolorbox{theorybox}[1]{
  colback=theorycolor!7, colframe=theorycolor,
  fonttitle=\bfseries\color{white}, colbacktitle=theorycolor,
  title={#1}, breakable, enhanced}

\newtcolorbox{practbox}[1]{
  colback=practcolor!8, colframe=practcolor,
  fonttitle=\bfseries\color{white}, colbacktitle=practcolor,
  title={#1}, breakable, enhanced}

\newtcolorbox{boardbox}[1]{
  colback=board!92, colframe=board!50!black,
  fonttitle=\bfseries\color{chalkyellow}, colbacktitle=board!80!black,
  title={#1}, breakable, enhanced, boxrule=2pt}

\newtcolorbox{nota}[1][]{
  colback=chalkyellow!50, colframe=orange!70,
  left=4pt, right=4pt, top=2pt, bottom=2pt,
  breakable, enhanced, sharp corners, #1}

\newtcolorbox{aviso}[1][]{
  colback=warncolor!8, colframe=warncolor,
  left=4pt, right=4pt, top=2pt, bottom=2pt,
  breakable, enhanced, sharp corners, #1}

%% ══════════════════════════════════════════════════════════════════
%% COMANDOS AUXILIARES
%% ══════════════════════════════════════════════════════════════════
\newcommand{\tiempo}[1]{%
  \vspace{4pt}
  \begin{tcolorbox}[
    colback=timercolor!12, colframe=timercolor,
    left=4pt, right=4pt, top=3pt, bottom=3pt,
    boxrule=1pt, sharp corners, breakable, enhanced]
  \textbf{\large\textcolor{timercolor}{\faIcon{clock}\quad #1}}
  \end{tcolorbox}\vspace{2pt}}

\newcommand{\tecla}[1]{\fbox{\small\texttt{#1}}}

%% ══════════════════════════════════════════════════════════════════
%% ESTILOS TIKZ GLOBALES
%% ══════════════════════════════════════════════════════════════════
\tikzset{
  obj/.style={
    rounded corners=7pt, minimum width=3.2cm, minimum height=1.0cm,
    align=center, font=\small\bfseries, draw, very thick},
  flecha/.style={-{Latex[length=3.5mm]}, very thick, darkblue},
  etiq/.style={font=\footnotesize\bfseries},
}

%% ══════════════════════════════════════════════════════════════════
\begin{document}
%% ══════════════════════════════════════════════════════════════════

%% ── PORTADA ──────────────────────────────────────────────────────
\begin{tcolorbox}[
  colback=accessred!5, colframe=accessred,
  boxrule=3pt, arc=4pt, left=12pt, right=12pt, top=10pt, bottom=10pt,
  breakable, enhanced]
\begin{center}
  {\color{accessred}\rule{\linewidth}{2pt}}\\[4pt]
  {\large\bfseries\color{accessred}
    Guión Docente — Sesiones 1 y 2\\[4pt]
    \large UT6: Microsoft Access}\\[6pt]
  {\normalsize\color{darkblue} Módulo: Aplicaciones Ofimáticas (AOF)}\\[4pt]
  {\color{accessred}\rule{\linewidth}{1pt}}\\[8pt]
  \begin{tabular}{>{\centering\arraybackslash}p{4cm}
                  >{\centering\arraybackslash}p{4cm}
                  >{\centering\arraybackslash}p{4cm}}
    \textbf{Sesión 1} & \textbf{Sesión 2} & \textbf{Total} \\[2pt]
    \small Conceptos básicos\\e interfaz &
    \small Crear tablas:\\Vista Diseño &
    \small 2 $\times$ 60\,min \\
  \end{tabular}
\end{center}
\end{tcolorbox}

\bigskip

\begin{tcolorbox}[colback=graylight, colframe=gray!40, sharp corners,
                  left=4pt, right=4pt, top=4pt, bottom=4pt,
                  breakable, enhanced]
\small\textbf{Leyenda:}\quad
\textcolor{timercolor}{\faIcon{clock}} Tiempo\quad
\textcolor{actcolor}{\faIcon{bolt}} Activador\quad
\textcolor{theorycolor}{\faIcon{chalkboard-teacher}} Teoría\quad
\textcolor{practcolor}{\faIcon{laptop}} Práctica\quad
\textcolor{board!60!black}{\faIcon{chalkboard}} Pizarra\quad
\textcolor{tipcolor}{\faIcon{lightbulb}} Consejo\quad
\textcolor{warncolor}{\faIcon{exclamation-triangle}} Atención
\end{tcolorbox}

\newpage

%% ════════════════════════════════════════════════════════════════
\section{\faIcon{database}\quad Sesión 1 · Conceptos básicos e interfaz}
%% ════════════════════════════════════════════════════════════════

\begin{tcolorbox}[colback=accessred!5, colframe=accessred, sharp corners,
                  breakable, enhanced, left=6pt, right=6pt]
\textbf{Objetivos:}~
Entender qué es una base de datos y para qué sirve \textbf{·}
Conocer los 6 objetos de Access \textbf{·}
Familiarizarse con la interfaz.
\end{tcolorbox}

%% ── 0‑5 min ──────────────────────────────────────────────────────
\tiempo{0 – 5 min \normalfont\normalsize | \faIcon{bolt} Activador}

\begin{actbox}{\faIcon{bolt}\quad Activador — «¿Para qué usamos datos?» (5 min)}
Pregunta: \textit{«¿Cuántas veces al día usáis una base de datos sin saberlo?»}

\medskip
Escribe en la pizarra lo que digan y dibuja el diagrama siguiente:

\begin{center}
\begin{tikzpicture}[node distance=0cm]
  \node[circle, fill=accessred!80, draw=accessred!60!black,
        text=white, font=\bfseries\small, minimum size=1.8cm,
        align=center] (c) {BBDD\\en tu\\vida};

  \node[rounded corners=4pt, fill=chalkyellow, draw=orange!70,
        font=\small\bfseries, minimum width=1.9cm, minimum height=0.65cm,
        align=center] (sp) at (90:3.2cm)  {Spotify};
  \node[rounded corners=4pt, fill=chalkyellow, draw=orange!70,
        font=\small\bfseries, minimum width=1.9cm, minimum height=0.65cm,
        align=center] (nf) at (45:3.2cm)  {Netflix};
  \node[rounded corners=4pt, fill=chalkyellow, draw=orange!70,
        font=\small\bfseries, minimum width=1.9cm, minimum height=0.65cm,
        align=center] (ba) at (0:3.2cm)   {Tu banco};
  \node[rounded corners=4pt, fill=chalkyellow, draw=orange!70,
        font=\small\bfseries, minimum width=1.9cm, minimum height=0.65cm,
        align=center] (md) at (315:3.2cm) {Médico};
  \node[rounded corners=4pt, fill=chalkyellow, draw=orange!70,
        font=\small\bfseries, minimum width=1.9cm, minimum height=0.65cm,
        align=center, text width=1.7cm] (ti) at (270:3.2cm) {Tienda online};
  \node[rounded corners=4pt, fill=chalkyellow, draw=orange!70,
        font=\small\bfseries, minimum width=1.9cm, minimum height=0.65cm,
        align=center] (ig) at (225:3.2cm) {Instagram};
  \node[rounded corners=4pt, fill=chalkyellow, draw=orange!70,
        font=\small\bfseries, minimum width=1.9cm, minimum height=0.65cm,
        align=center, text width=1.7cm] (co) at (180:3.2cm) {Notas cole};
  \node[rounded corners=4pt, fill=chalkyellow, draw=orange!70,
        font=\small\bfseries, minimum width=1.9cm, minimum height=0.65cm,
        align=center, text width=1.7cm] (gm) at (135:3.2cm) {Google Maps};

  \foreach \n in {sp,nf,ba,md,ti,ig,co,gm}
    \draw[->, gray!50, thick] (c) -- (\n);
\end{tikzpicture}
\end{center}

\textit{«Todo esto es una base de datos. Hoy vais a aprender a construir una.»}
\end{actbox}

%% ── 5‑20 min ─────────────────────────────────────────────────────
\tiempo{5 – 20 min \normalfont\normalsize | \faIcon{chalkboard-teacher} Teoría: conceptos básicos}

\begin{theorybox}{\faIcon{chalkboard-teacher}\quad Conceptos clave (15 min)}
\begin{itemize}[leftmargin=1.5em]
  \item \textbf{Base de datos}: conjunto de datos organizados para un fin.
  \item \textbf{SGBD}: software que gestiona esos datos. Access es un SGBD.
  \item Modelo relacional: datos en \textbf{tablas} relacionadas. Extensión: \texttt{.accdb}
  \item \textbf{Clave primaria (PK)}: identifica de forma única cada fila.
  \item \textbf{Clave foránea (FK)}: apunta a la PK de otra tabla.
  \item \textbf{NULL}: ausencia de dato (distinto de 0 o cadena vacía).
\end{itemize}

\medskip
Dibuja en la pizarra las \textbf{Pizarras 1, 2 y 3} que siguen.
\end{theorybox}

%% ── PIZARRA 1 ─────────────────────────────────────────────────────
\begin{boardbox}{\faIcon{chalkboard}\quad PIZARRA 1 — Anatomía de una tabla}

% ── Tabla con colores ──
\begin{center}
\renewcommand{\arraystretch}{1.3}
\begin{tabular}{|>{\columncolor{chalkyellow!60}}c|c|c|c|}
\hline
\rowcolor{darkblue!80}
\color{white}\bfseries\faIcon{key}~CodCliente &
\color{white}\bfseries Nombre &
\color{white}\bfseries Apellidos &
\color{white}\bfseries Ciudad \\
\hline
\rowcolor{chalkwhite} C001 & Ana   & García López   & Valencia \\
\rowcolor{chalk}      C002 & Pedro & Martínez Gil   & Madrid   \\
\rowcolor{chalkwhite} C003 & María & García López   & Sevilla  \\
\rowcolor{chalk}      C004 & Luisa & Fernández Díaz & Valencia \\
\hline
\end{tabular}
\end{center}

\begin{tikzpicture}[overlay, remember picture]
\end{tikzpicture}

\vspace{4pt}
% ── Anotaciones ──
\begin{center}
\begin{tikzpicture}
  \node[draw=darkblue, fill=lightblue!40, rounded corners=5pt,
        minimum width=3cm, font=\small, align=center, inner sep=5pt]
    (campo) at (0,0) {\textbf{\textcolor{darkblue}{CAMPO}}\\(columna)\\{\footnotesize Mismo tipo\\en toda la columna}};

  \node[draw=accessred, fill=chalkred!40, rounded corners=5pt,
        minimum width=3cm, font=\small, align=center, inner sep=5pt,
        right=2cm of campo]
    (registro) {\textbf{\textcolor{accessred}{REGISTRO}}\\(fila)\\{\footnotesize Todos los datos\\de un elemento}};

  \node[draw=orange!80, fill=chalkyellow!70, rounded corners=5pt,
        minimum width=3cm, font=\small, align=center, inner sep=5pt,
        right=2cm of registro]
    (pk) {\textbf{\textcolor{orange!80!black}{\faIcon{key}~CLAVE PK}}\\{\footnotesize Único · No nulo\\No repetible}};
\end{tikzpicture}
\end{center}

\begin{nota}
\small\textbf{Pregunta al grupo:} ¿Podría «Apellidos» ser clave primaria?
\textit{No: C001 y C003 tienen los mismos apellidos.}
¿Y «Nombre»? \textit{No: puede haber dos «María».}
\end{nota}
\end{boardbox}

%% ── PIZARRA 2 ─────────────────────────────────────────────────────
\begin{boardbox}{\faIcon{chalkboard}\quad PIZARRA 2 — Relación Clientes $\rightarrow$ Pedidos}

\begin{center}
\begin{tikzpicture}

  %% ── Tabla CLIENTES (izquierda) ──
  \node[fill=darkblue!75, text=white, font=\bfseries\small,
        minimum width=5cm, minimum height=0.65cm]
    (hC) at (0,0) {CLIENTES};

  % Fila PK
  \node[fill=chalkyellow!70, draw=orange!70, minimum width=5cm,
        minimum height=0.55cm, font=\small\bfseries] (c1) at (0,-0.6)
    {\faIcon{key}~CodCliente \quad {\footnotesize\normalfont Texto corto}};

  % Filas normales
  \node[fill=chalkwhite, draw=gray!30, minimum width=5cm,
        minimum height=0.55cm, font=\small] (c2) at (0,-1.15)
    {Nombre \hfill {\footnotesize Texto corto}};
  \node[fill=chalk, draw=gray!30, minimum width=5cm,
        minimum height=0.55cm, font=\small] (c3) at (0,-1.70)
    {Apellidos \hfill {\footnotesize Texto corto}};
  \node[fill=chalkwhite, draw=gray!30, minimum width=5cm,
        minimum height=0.55cm, font=\small] (c4) at (0,-2.25)
    {Ciudad \hfill {\footnotesize Texto corto}};
  \node[fill=chalk, draw=gray!30, minimum width=5cm,
        minimum height=0.55cm, font=\small] (c5) at (0,-2.80)
    {FechaNacimiento \hfill {\footnotesize Fecha/Hora}};

  %% ── Tabla PEDIDOS (derecha) ──
  \node[fill=darkgreen!70, text=white, font=\bfseries\small,
        minimum width=5cm, minimum height=0.65cm]
    (hP) at (7.5,0) {PEDIDOS};

  % Fila PK
  \node[fill=chalkyellow!70, draw=orange!70, minimum width=5cm,
        minimum height=0.55cm, font=\small\bfseries] (p1) at (7.5,-0.6)
    {\faIcon{key}~CodPedido \quad {\footnotesize\normalfont Autonumeración}};

  % Fila FK
  \node[fill=chalkred!50, draw=accessred!60, minimum width=5cm,
        minimum height=0.55cm, font=\small\bfseries,
        text=accessred!80!black] (p2) at (7.5,-1.15)
    {\faIcon{link}~CodCliente \quad {\footnotesize\normalfont Texto corto}};

  % Filas normales
  \node[fill=chalk, draw=gray!30, minimum width=5cm,
        minimum height=0.55cm, font=\small] (p3) at (7.5,-1.70)
    {FechaPedido \hfill {\footnotesize Fecha/Hora}};
  \node[fill=chalkwhite, draw=gray!30, minimum width=5cm,
        minimum height=0.55cm, font=\small] (p4) at (7.5,-2.25)
    {Enviado \hfill {\footnotesize Sí/No}};

  %% ── Línea de relación ──
  \draw[very thick, -{Latex[length=4mm]}, darkblue!70]
    (c1.east) .. controls (3.5,-0.6) and (4.0,-1.15) .. (p2.west);

  \node[font=\bfseries\small, text=darkblue] at (2.8,-0.3) {1};
  \node[font=\bfseries\small, text=darkblue] at (5.0,-1.4) {N};

  %% ── Leyenda ──
  \node[fill=chalkyellow!70, draw=orange!70, minimum width=0.5cm,
        minimum height=0.4cm] (lk) at (0,-3.6) {};
  \node[font=\footnotesize, right=0.15cm of lk] {\faIcon{key} Clave primaria (PK)};

  \node[fill=chalkred!50, draw=accessred!60, minimum width=0.5cm,
        minimum height=0.4cm] (lfk) at (4.5,-3.6) {};
  \node[font=\footnotesize, right=0.15cm of lfk]
    {\faIcon{link} Clave foránea (FK)};

\end{tikzpicture}
\end{center}

\begin{nota}
\small\textbf{Explica:}
CodCliente en PEDIDOS es FK porque «apunta» a un cliente de la tabla CLIENTES.
Un cliente puede tener muchos pedidos $\Rightarrow$ relación \textbf{1 a N (uno a varios)}.
\end{nota}
\end{boardbox}

%% ── PIZARRA 3: NULL vs 0 ──────────────────────────────────────────
\begin{boardbox}{\faIcon{chalkboard}\quad PIZARRA 3 — NULL vs. 0}

\begin{center}
\begin{tikzpicture}
  \node[draw=accessred, fill=chalkred!30, rounded corners=8pt,
        minimum width=5cm, minimum height=2cm,
        align=center, font=\normalsize, inner sep=8pt] (n) at (0,0)
    {\textbf{\textcolor{accessred}{NULL}}\\[4pt]
     «No tengo ese dato»\\[2pt]
     \textit{\small «No sé su edad»}};

  \node[draw=darkgreen!70, fill=lightgreen!50, rounded corners=8pt,
        minimum width=5cm, minimum height=2cm,
        align=center, font=\normalsize, inner sep=8pt,
        right=2.5cm of n] (z)
    {\textbf{\textcolor{darkgreen}{0}}\\[4pt]
     «El dato existe y vale cero»\\[2pt]
     \textit{\small «Tiene 0 hijos»}};

  \node[font=\Huge\bfseries, text=accessred]
    at ($(n.east)!0.5!(z.west)$) {$\neq$};
\end{tikzpicture}
\end{center}
\end{boardbox}

%% ── PIZARRA 4: Los 6 objetos ──────────────────────────────────────
\begin{boardbox}{\faIcon{chalkboard}\quad PIZARRA 4 — Los 6 objetos de Access}

\begin{center}
\begin{tikzpicture}
  %% Centro
  \node[circle, fill=accessred!80, draw=accessred!60!black,
        text=white, font=\bfseries\small, minimum size=1.8cm,
        align=center] (c) {ACCESS\\DB};

  %% Los 6 objetos
  \node[obj, fill=lightblue!60,   draw=darkblue,     above=2.0cm of c]
    (tabla) {\faIcon{table}\\Tabla\\{\footnotesize\normalfont Almacena datos}};

  \node[obj, fill=lightgreen!60,  draw=darkgreen!70,
        above right=1.2cm and 2.6cm of c]
    (consulta) {\faIcon{search}\\Consulta\\{\footnotesize\normalfont Filtra y calcula}};

  \node[obj, fill=chalkyellow!70, draw=orange!70,    right=3.4cm of c]
    (form) {\faIcon{window-maximize}\\Formulario\\{\footnotesize\normalfont Interfaz visual}};

  \node[obj, fill=chalkred!40,    draw=accessred!60,
        below right=1.2cm and 2.6cm of c]
    (informe) {\faIcon{print}\\Informe\\{\footnotesize\normalfont Para imprimir}};

  \node[obj, fill=gray!20,        draw=gray!60,      below=2.0cm of c]
    (macro) {\faIcon{robot}\\Macro\\{\footnotesize\normalfont Automatiza acciones}};

  \node[obj, fill=gray!15,        draw=gray!50,
        below left=1.2cm and 2.6cm of c]
    (modulo) {\faIcon{code}\\Módulo\\{\footnotesize\normalfont Código VBA}};

  %% Flechas al centro
  \foreach \n in {tabla,consulta,form,informe,macro,modulo}
    \draw[<->, thick, gray!50] (c) -- (\n);

  %% Flechas de dependencia
  \draw[->, thick, darkblue!60, dashed, bend left=25]
    (tabla) to node[above, font=\footnotesize, text=darkblue]{base} (consulta);
  \draw[->, thick, darkblue!60, dashed, bend left=15]
    (tabla) to (form);
  \draw[->, thick, darkblue!60, dashed, bend right=15]
    (tabla) to (informe);
\end{tikzpicture}
\end{center}

\begin{nota}
\small Las \textbf{tablas} son la base de todo.
Consultas, formularios e informes se construyen \emph{siempre} sobre tablas.
\textbf{En este módulo}: tabla, consulta y formulario.
\end{nota}
\end{boardbox}

%% ── 20‑35 min ────────────────────────────────────────────────────
\tiempo{20 – 35 min \normalfont\normalsize | \faIcon{chalkboard-teacher} Interfaz de Access}

\begin{theorybox}{\faIcon{chalkboard-teacher}\quad Demo en proyector — Interfaz (15 min)}

Recorre en directo con el proyector y señala cada elemento:
\begin{enumerate}[leftmargin=1.8em]
  \item Pantalla inicial: crear en blanco, plantillas, recientes.
  \item Abrir \texttt{TiendaInformatica.accdb}.
  \item Señalar cada elemento del entorno (ver \textbf{Pizarra 5}).
  \item Demo de la cinta: pestaña $\rightarrow$ grupo $\rightarrow$ botón.
        Notación: \texttt{INICIO > Portapapeles > Pegar}.
  \item Vista Backstage: pestaña \textbf{Archivo}.
\end{enumerate}

\begin{tcolorbox}[colback=actcolor!10, colframe=actcolor, sharp corners,
                  left=4pt, right=4pt, top=2pt, bottom=2pt,
                  breakable, enhanced]
\faIcon{bolt}\;\textbf{Actividad rápida (2 min):}
¿Cuántas formas hay de guardar en Access?
El primero en encontrar 3 formas distintas gana un punto.\\
\textit{(Ctrl+G · botón Guardar en barra de acceso rápido · Archivo > Guardar)}
\end{tcolorbox}
\end{theorybox}

%% ── PIZARRA 5 ─────────────────────────────────────────────────────
\begin{boardbox}{\faIcon{chalkboard}\quad PIZARRA 5 — Elementos de la interfaz de Access}

\begin{center}
\begin{tikzpicture}[scale=0.93, font=\small]

  %% Marco general de la ventana
  \fill[gray!12] (0,0) rectangle (14,9);
  \draw[gray!50, very thick] (0,0) rectangle (14,9);

  %% 1. Barra de título
  \fill[gray!45] (0,8.4) rectangle (14,9);
  \node[font=\footnotesize\bfseries, text=white] at (7,8.7)
    {TiendaInformatica : Base de datos (Access)};
  \draw[->, red!70!black, thick] (7,9.45) -- (7,9.05)
    node[above, at start, font=\footnotesize\bfseries, text=red!70!black]
    {\textbf{1.} Barra de título};

  %% 2. Barra de acceso rápido
  \fill[gray!38] (0,8.0) rectangle (2.4,8.4);
  \node[font=\tiny, text=white] at (1.2,8.2)
    {\faIcon{save}\quad\faIcon{undo}\quad\faIcon{redo}};
  \draw[->, darkblue, thick] (-0.9,8.2) -- (-0.05,8.2)
    node[left, at start, font=\footnotesize\bfseries, text=darkblue, align=right]
    {\textbf{2.} Barra acceso\\rápido};

  %% 3. Cinta de opciones
  \fill[accessred!12] (0,7.1) rectangle (14,8.0);
  % pestañas
  \foreach \xp/\nm in {0.55/ARCHIVO, 1.65/INICIO, 2.65/CREAR,
                        3.75/{DAT.EXT.}, 5.1/{H.BBDD}}{
    \fill[white] (\xp-0.5,7.65) rectangle (\xp+0.55,8.0);
    \node[font=\tiny\bfseries] at (\xp+0.02,7.82) {\nm};
  }
  % grupos en cinta
  \foreach \xg/\gn in {0.9/{Portapapeles}, 2.2/{Vistas},
                        3.5/{Ordenar}, 4.9/{Registros}, 6.2/{Buscar}}{
    \draw[gray!40] (\xg,7.14) rectangle (\xg+1.1,7.62);
    \node[font=\tiny, text=gray!60] at (\xg+0.55,7.2) {\gn};
    \foreach \bx in {0.15,0.45,0.75}{
      \fill[lightblue!70] (\xg+\bx,7.32) rectangle (\xg+\bx+0.24,7.56);
    }
  }
  \draw[->, accessred, thick] (14.8,7.55) -- (14.05,7.55)
    node[right, at start, font=\footnotesize\bfseries, text=accessred, align=left]
    {\textbf{3.} Cinta de opciones\\(pestañas $\to$ grupos)};

  %% 4. Panel de navegación
  \fill[white] (0,0.4) rectangle (2.8,7.1);
  \draw[gray!40] (0,0.4) rectangle (2.8,7.1);
  \fill[darkblue!70] (0,6.7) rectangle (2.8,7.1);
  \node[font=\tiny\bfseries, text=white] at (1.4,6.9)
    {Todos los objetos Access};
  \foreach \yp/\ic/\lb in {
    6.4/{\faIcon{table}}/{Tablas},
    5.8/{\faIcon{search}}/{Consultas},
    5.2/{\faIcon{window-maximize}}/{Formularios},
    4.6/{\faIcon{print}}/{Informes}}{
    \fill[gray!10] (0.1,\yp-0.22) rectangle (2.7,\yp+0.22);
    \node[font=\tiny\bfseries, text=darkblue] at (0.5,\yp) {\ic};
    \node[font=\tiny] at (1.6,\yp) {\lb};
  }
  \draw[->, darkgreen!70, thick] (-0.9,5.5) -- (-0.05,5.5)
    node[left, at start, font=\footnotesize\bfseries, text=darkgreen!70, align=right]
    {\textbf{4.} Panel de\\navegación};

  %% 5. Área de trabajo
  \fill[white] (2.9,0.4) rectangle (13.9,7.0);
  \draw[gray!35] (2.9,0.4) rectangle (13.9,7.0);
  % pestaña de tabla
  \fill[lightblue!35] (2.9,6.6) rectangle (5.0,7.0);
  \node[font=\tiny\bfseries] at (3.95,6.8) {\faIcon{table}~Clientes};
  % cabecera tabla
  \fill[darkblue!70] (3.0,6.15) rectangle (13.8,6.58);
  \foreach \xh/\hd in {4.0/{CodCliente}, 6.3/{Nombre},
                        8.7/{Apellidos}, 11.2/{Ciudad}}{
    \node[font=\tiny\bfseries, text=white] at (\xh,6.36) {\hd};
  }
  % filas de datos
  \foreach \yi/\bgc/\cod/\nm/\ap/\ci in {
    1/chalkwhite/C001/Ana/García/Valencia,
    2/chalk/C002/Pedro/Martínez/Madrid,
    3/chalkwhite/C003/María/Fernández/Sevilla}{
    \fill[\bgc] (3.0,{6.15-\yi*0.45}) rectangle (13.8,{6.15-\yi*0.45+0.43});
    \node[font=\tiny] at (4.0,{6.15-\yi*0.45+0.215}) {\texttt{\cod}};
    \node[font=\tiny] at (6.3,{6.15-\yi*0.45+0.215}) {\nm};
    \node[font=\tiny] at (8.7,{6.15-\yi*0.45+0.215}) {\ap};
    \node[font=\tiny] at (11.2,{6.15-\yi*0.45+0.215}) {\ci};
  }
  \draw[->, gray!55, thick] (14.8,4.5) -- (13.95,4.5)
    node[right, at start, font=\footnotesize\bfseries, text=gray!55, align=left]
    {\textbf{5.} Área de\\trabajo};

  %% 6. Barra de estado
  \fill[gray!35] (0,0) rectangle (14,0.4);
  \node[font=\tiny, text=white] at (3.0,0.2) {Vista: Hoja de datos};
  \foreach \xb in {11.4,12.0,12.6,13.2}{
    \fill[gray!60] (\xb-0.2,0.08) rectangle (\xb+0.2,0.32);
  }
  \draw[->, timercolor, thick] (-0.9,0.2) -- (-0.05,0.2)
    node[left, at start, font=\footnotesize\bfseries, text=timercolor, align=right]
    {\textbf{6.} Barra de\\estado};

\end{tikzpicture}
\end{center}
\end{boardbox}

%% ── 35‑55 min ────────────────────────────────────────────────────
\tiempo{35 – 55 min \normalfont\normalsize | \faIcon{laptop} Práctica}

\begin{practbox}{\faIcon{laptop}\quad Práctica (20 min) — Explorar la interfaz}
\begin{enumerate}[leftmargin=1.8em]
  \item Abrir Access y crear \texttt{TiendaInformatica.accdb} en la carpeta personal.
  \item Identificar en pantalla cada uno de los 6 elementos de la Pizarra 5.
  \item Contraer y expandir el panel de navegación y la cinta de opciones.
  \item Ejercicio 1.1 del cuaderno (vocabulario).
  \item Ejercicio 1.2 del cuaderno (analizar la tabla de clientes).
\end{enumerate}

\begin{aviso}
\faIcon{exclamation-triangle}\;
\textbf{Circula por las mesas.}
Los elementos más difíciles de localizar son la barra de acceso rápido y la barra de estado.
\end{aviso}
\end{practbox}

%% ── 55‑60 min ────────────────────────────────────────────────────
\tiempo{55 – 60 min \normalfont\normalsize | Cierre y debate}

\begin{actbox}{\faIcon{comments}\quad Cierre — Debate rápido (5 min)}
\textit{«¿Qué ventajas tiene Access frente a Excel para gestionar datos de una tienda?»}

\medskip
\begin{center}
\begin{tabular}{p{5.5cm} p{5.5cm}}
\hline
\rowcolor{lightblue!40}
\textbf{\textcolor{darkblue}{Access \faIcon{database}}} &
\textbf{\textcolor{darkgreen}{Excel \faIcon{table}}} \\
\hline
$\bullet$ Múltiples tablas relacionadas &
$\bullet$ Más fácil de aprender \\
\rowcolor{gray!8}
$\bullet$ Sin duplicar datos &
$\bullet$ Gráficos integrados \\
$\bullet$ Consultas potentes &
$\bullet$ Cálculos ad-hoc \\
\rowcolor{gray!8}
$\bullet$ Formularios para usuarios &
$\bullet$ Todo visible a la vez \\
$\bullet$ Integridad de datos &
$\bullet$ Más extendido en empresas \\
\hline
\end{tabular}
\end{center}

\medskip
\textbf{Anticipa la sesión 2:}
\textit{«La próxima sesión creamos nuestras primeras tablas y veréis por qué un teléfono NUNCA debe ser de tipo Número.»}
\end{actbox}

%% ════════════════════════════════════════════════════════════════
\newpage
\section{\faIcon{table}\quad Sesión 2 · Crear tablas: Vista Diseño}
%% ════════════════════════════════════════════════════════════════

\begin{tcolorbox}[colback=accessred!5, colframe=accessred, sharp corners,
                  breakable, enhanced, left=6pt, right=6pt]
\textbf{Objetivos:}~
Crear tablas en Vista Diseño \textbf{·}
Elegir correctamente el tipo de dato \textbf{·}
Asignar clave principal y guardar la tabla.
\end{tcolorbox}

%% ── 0‑5 min ──────────────────────────────────────────────────────
\tiempo{0 – 5 min \normalfont\normalsize | \faIcon{bolt} Repaso exprés}

\begin{actbox}{\faIcon{bolt}\quad Repaso (5 min) — Preguntas orales al azar}
\begin{enumerate}[leftmargin=1.8em]
  \item ¿Qué es un campo? ¿Y un registro?
  \item ¿Puede haber dos clientes con el mismo CodCliente? ¿Por qué?
  \item ¿Cuántos objetos tiene una base de datos Access? Nómbralos.
  \item ¿Qué diferencia hay entre NULL y el valor 0?
\end{enumerate}
\textit{Si alguien responde muy bien, pídele que pase a la pizarra a dibujar una tabla rápida.}
\end{actbox}

%% ── 5‑20 min ─────────────────────────────────────────────────────
\tiempo{5 – 20 min \normalfont\normalsize | \faIcon{chalkboard-teacher} Tipos de datos}

\begin{theorybox}{\faIcon{chalkboard-teacher}\quad Teoría (15 min) — Tipos de datos}

\textbf{Regla de oro:} el tipo de dato determina qué valores puede almacenar el campo y cuánto espacio ocupa.

\begin{tcolorbox}[colback=actcolor!10, colframe=actcolor, sharp corners,
                  left=4pt, right=4pt, top=2pt, bottom=2pt,
                  breakable, enhanced]
\faIcon{bolt}\;\textbf{Juego rápido (3 min):}
dicta 10 campos, los alumnos apuntan el tipo correcto en su cuaderno.
Corrección oral a continuación.
\end{tcolorbox}

Dibuja en la pizarra la \textbf{Pizarra 6} que sigue.
\end{theorybox}

%% ── PIZARRA 6 ─────────────────────────────────────────────────────
\begin{boardbox}{\faIcon{chalkboard}\quad PIZARRA 6 — Tipos de datos de Access}

\begin{center}
\renewcommand{\arraystretch}{1.35}
\begin{tabularx}{\linewidth}{
    >{\bfseries\color{darkblue}}p{3.0cm}
    p{2.8cm}
    X
    >{\itshape\color{gray!70!black}\small}p{3.2cm}}
\toprule
\normalfont\bfseries\color{white}\cellcolor{darkblue!80} Tipo &
\cellcolor{darkblue!80}\color{white}\bfseries Tamaño &
\cellcolor{darkblue!80}\color{white}\bfseries Úsalo cuando\ldots &
\cellcolor{darkblue!80}\color{white}\bfseries Ejemplo \\
\midrule
\rowcolor{chalkwhite}
Texto corto & hasta 255 car. & texto, códigos, teléfonos, CP &
  Nombre, DNI, CP \\
\rowcolor{chalk}
Texto largo & hasta 1\,GB & descripciones largas &
  Observaciones \\
\rowcolor{chalkwhite}
Número & Byte/Entero/Doble & cálculos matemáticos &
  Stock, Cantidad \\
\rowcolor{chalk}
Fecha / Hora & 8 bytes & fechas y horas &
  FechaNacimiento \\
\rowcolor{chalkwhite}
Moneda & 8 bytes & importes, precios &
  PrecioUnidad \\
\rowcolor{chalk}
Autonumeración & 4 bytes & código auto (ideal para PK) &
  CodPedido \\
\rowcolor{chalkwhite}
Sí / No & 1 bit & datos booleanos &
  Enviado, Activo \\
\rowcolor{chalk}
Datos adjuntos & variable & archivos e imágenes &
  FotoProducto \\
\rowcolor{chalkwhite}
Hipervínculo & variable & URL, ruta de archivo &
  WebEmpresa \\
\rowcolor{chalk}
Calculado & según expr. & resultado de expresión &
  Precio $\times$ Cantidad \\
\bottomrule
\end{tabularx}
\end{center}

\begin{aviso}
\faIcon{exclamation-triangle}\;\textbf{Error clásico:}
El teléfono y el código postal son \textbf{Texto corto}, nunca Número.
El \texttt{0} inicial desaparece si usas Número (\texttt{698765432} $\to$ ¡pierdes el prefijo!).
Además, no se hacen cálculos con ellos.
\end{aviso}
\end{boardbox}

%% ── 20‑30 min ────────────────────────────────────────────────────
\tiempo{20 – 30 min \normalfont\normalsize | \faIcon{chalkboard-teacher} Demo en proyector}

\begin{theorybox}{\faIcon{chalkboard-teacher}\quad Demo — Vista Diseño paso a paso (10 min)}
\begin{enumerate}[leftmargin=1.8em]
  \item \texttt{Crear > Tablas > Vista Diseño}
  \item Escribir nombre del campo, elegir tipo, añadir descripción (opcional).
  \item Repetir para cada campo.
  \item Asignar clave principal: clic en el campo
        $\rightarrow$ \texttt{Diseño > Clave principal}. Aparece la llave \faIcon{key}.
  \item Guardar: \tecla{Ctrl+G} $\rightarrow$ nombre $\rightarrow$ Aceptar.
  \item Abrir en Vista Hoja de datos e introducir 2 registros.
\end{enumerate}

Dibuja en la pizarra la \textbf{Pizarra 7}.
\end{theorybox}

%% ── PIZARRA 7 ─────────────────────────────────────────────────────
\begin{boardbox}{\faIcon{chalkboard}\quad PIZARRA 7 — Vista Diseño anotada}

\begin{center}
\renewcommand{\arraystretch}{1.3}
\begin{tabular}{|>{\columncolor{chalkyellow!60}\bfseries\small\color{orange!80!black}}p{3.2cm}
                |>{\small\color{darkblue}\bfseries}p{2.8cm}
                |>{\small\itshape\color{gray!60!black}}p{5.5cm}|}
\hline
\rowcolor{darkblue!80}
\color{white}\bfseries\normalsize\normalfont Nombre del campo &
\color{white}\bfseries\normalsize\normalfont Tipo de datos &
\color{white}\bfseries\normalsize\normalfont Descripción (opcional) \\
\hline
\faIcon{key}~CodCliente   & Texto corto   & Código único del cliente \\
\hline
\rowcolor{chalk}
Nombre           & Texto corto   & Nombre de pila \\
\hline
Apellidos        & Texto corto   & Apellidos completos \\
\hline
\rowcolor{chalk}
Ciudad           & Texto corto   & Ciudad de residencia \\
\hline
FechaNacimiento  & Fecha / Hora  & Fecha de nacimiento \\
\hline
\rowcolor{chalk}
Activo           & Sí / No       & ¿Sigue siendo cliente activo? \\
\hline
\end{tabular}
\end{center}

\vspace{8pt}
\begin{center}
\begin{tikzpicture}[node distance=0.4cm]
  \node[draw=orange!80, fill=chalkyellow!60, rounded corners=5pt,
        font=\small, align=center, inner sep=6pt, minimum width=3cm] (pk)
    {\faIcon{key}~\textbf{Clave primaria}\\{\footnotesize Fila con fondo amarillo}\\{\footnotesize Único $\cdot$ No nulo}};

  \node[draw=darkblue, fill=lightblue!40, rounded corners=5pt,
        font=\small, align=center, inner sep=6pt, minimum width=3.5cm,
        right=1.5cm of pk] (prop)
    {\textbf{Propiedades del campo}\\{\footnotesize (zona inferior)}\\{\footnotesize Tamaño, Formato, Máscara,}\\{\footnotesize Valor predeterminado\ldots}};

  \node[draw=darkgreen!70, fill=lightgreen!40, rounded corners=5pt,
        font=\small, align=center, inner sep=6pt, minimum width=3cm,
        right=1.5cm of prop] (tipo)
    {\textbf{Tipo de datos}\\{\footnotesize Cambia las propiedades}\\{\footnotesize disponibles en el panel}\\{\footnotesize inferior}};
\end{tikzpicture}
\end{center}
\end{boardbox}

%% ── PIZARRA 8: Flujo crear tabla ──────────────────────────────────
\begin{boardbox}{\faIcon{chalkboard}\quad PIZARRA 8 — Flujo completo: crear una tabla}

\begin{center}
\begin{tikzpicture}[node distance=0.6cm and 0.7cm]
  \node[obj, fill=darkblue!70, draw=darkblue!50!black, text=white,
        minimum width=2.8cm] (s1)
    {\faIcon{plus-circle}\\Crear\\Vista Diseño};

  \node[obj, fill=darkblue!50, draw=darkblue!40!black, text=white,
        minimum width=2.8cm, right=0.8cm of s1] (s2)
    {\faIcon{keyboard}\\Definir campos\\(nombre + tipo)};

  \node[obj, fill=darkgreen!60, draw=darkgreen!40!black, text=white,
        minimum width=2.8cm, right=0.8cm of s2] (s3)
    {\faIcon{key}\\Clave\\principal};

  \node[obj, fill=tipcolor, draw=tipcolor!50!black, text=white,
        minimum width=2.8cm, right=0.8cm of s3] (s4)
    {\faIcon{save}\\Guardar\\\tecla{Ctrl+G}};

  \node[obj, fill=accessred!70, draw=accessred!50!black, text=white,
        minimum width=2.8cm, right=0.8cm of s4] (s5)
    {\faIcon{table}\\Vista Hoja\\de datos};

  \foreach \a/\b in {s1/s2, s2/s3, s3/s4, s4/s5}
    \draw[flecha] (\a) -- (\b);

  \node[font=\footnotesize, text=gray, below=0.3cm of s2, align=center]
    {Un campo\\por fila};
  \node[font=\footnotesize, text=orange!80!black, below=0.3cm of s3, align=center]
    {Aparece\\la llave \faIcon{key}};
  \node[font=\footnotesize, text=tipcolor, below=0.3cm of s4, align=center]
    {Nombre\\de la tabla};
  \node[font=\footnotesize, text=accessred, below=0.3cm of s5, align=center]
    {Introducir\\datos};
\end{tikzpicture}
\end{center}
\end{boardbox}

%% ── 30‑57 min ────────────────────────────────────────────────────
\tiempo{30 – 57 min \normalfont\normalsize | \faIcon{laptop} Práctica guiada}

\begin{practbox}{\faIcon{laptop}\quad Práctica (27 min) — Crear las tablas de TiendaInformatica}

\textbf{Parte 1 (12 min): Tabla Clientes}
\begin{itemize}[leftmargin=1.8em]
  \item Seguir la ficha del cuaderno de ejercicios (sesión 2).
  \item Asignar PK al campo CodCliente.
  \item Introducir 5 registros en Vista Hoja de datos.
\end{itemize}

\textbf{Parte 2 (10 min): Tabla Productos}
\begin{itemize}[leftmargin=1.8em]
  \item CodProducto (Autonumeración, PK), Nombre (Texto, 100), Categoria (Texto, 30),
        PrecioUnidad (Moneda), Stock (Número entero).
  \item Introducir 5 productos.
\end{itemize}

\textbf{Parte 3 (5 min): «Caza el error»} — Por parejas:

\begin{center}
\begin{tcolorbox}[colback=chalkred!20, colframe=accessred, sharp corners,
                  left=6pt, right=6pt, breakable, enhanced]
\small Mostrar en el proyector esta tabla mal diseñada.
En 3 minutos encuentran todos los errores. La pareja que más encuentre y los justifique, gana.

\medskip
\renewcommand{\arraystretch}{1.2}
\begin{tabular}{|l|l|l|}
\hline
\rowcolor{darkblue!60}
\color{white}\bfseries Campo & \color{white}\bfseries Tipo asignado & \color{white}\bfseries ¿Error? \\
\hline
Teléfono & Número & \rule{3cm}{0.4pt} \\
\hline
\rowcolor{chalk}
Código postal & Número & \rule{3cm}{0.4pt} \\
\hline
Precio & Texto corto & \rule{3cm}{0.4pt} \\
\hline
\rowcolor{chalk}
FechaNacimiento & Texto corto & \rule{3cm}{0.4pt} \\
\hline
Descripción larga & Texto corto (50) & \rule{3cm}{0.4pt} \\
\hline
\rowcolor{chalk}
Activo (sí o no) & Número & \rule{3cm}{0.4pt} \\
\hline
\end{tabular}
\end{tcolorbox}
\end{center}

\begin{aviso}
\faIcon{exclamation-triangle}\;\textbf{Errores comunes a vigilar durante la práctica:}
\begin{itemize}[leftmargin=1.5em, itemsep=1pt]
  \item Usar Número para el teléfono o el código postal.
  \item No asignar clave principal antes de guardar.
  \item Confundir Moneda con Número para precios.
  \item Texto corto para campos de más de 255 caracteres.
\end{itemize}
\end{aviso}
\end{practbox}

%% ── 57‑60 min ────────────────────────────────────────────────────
\tiempo{57 – 60 min \normalfont\normalsize | Cierre}

\begin{actbox}{\faIcon{comments}\quad Cierre (3 min)}
\textit{«¿Por qué el tipo de dato de un campo es tan importante?
¿Qué podría pasar si guardamos el precio de un producto como Texto corto?»}

\smallskip
\textbf{Respuesta esperada:} Access lo tratará como texto — no podrá sumar totales,
calcular medias ni aplicar descuentos.

\medskip
\textbf{Anticipa la sesión 3:}
\textit{«La próxima sesión añadiremos propiedades a los campos: validaciones,
valores por defecto, máscaras\ldots{} Vuestras tablas se volverán a prueba de errores.»}
\end{actbox}

%% ── NOTAS FINALES ────────────────────────────────────────────────
\vspace{6pt}
\begin{tcolorbox}[colback=graylight, colframe=gray!40,
                  title=\textbf{Notas y recursos para el docente},
                  colbacktitle=gray!50, fonttitle=\bfseries\color{white},
                  breakable, enhanced]
\small
\begin{itemize}[leftmargin=1.5em, itemsep=2pt]
  \item \textbf{Archivo de práctica:} preparar \texttt{TiendaInformatica\_plantilla.accdb}
        con las tablas vacías para alumnos que se queden atrás.
  \item \textbf{Referencia libre:} \texttt{aulaclic.es/access-2019}
        — todo el contenido en español.
  \item \textbf{Si el tiempo aprieta:} omitir «Caza el error» y dejarlo como tarea.
  \item \textbf{Si el tiempo sobra:} reto extra del cuaderno (campo calculado Edad).
  \item \textbf{Material imprimible:} las pizarras pueden imprimirse en A3
        y dejarse en el aula como referencia visual permanente.
\end{itemize}
\end{tcolorbox}

\vfill
\begin{center}
{\color{accessred}\rule{0.4\linewidth}{1pt}}\\[4pt]
{\small\color{gray} Guión docente UT6 · Sesiones 1 y 2 · Microsoft Access · AOF}
\end{center}

\end{document}
